\begin{helper}
Let \(G\) be a finite \(\Delta\)-regular graph, let \(0<\gamma\le\Delta\),
and let \(\nu\) be a law on activation--priority pairs \((A,\pi)\).  Assume
that \(\nu\) has total mass one and that, for every \(A_0\subseteq V(G)\),
\[
\nu(A=A_0)=
\left(\frac{\gamma}{\Delta}\right)^{|A_0|}
\left(1-\frac{\gamma}{\Delta}\right)^{|V(G)|-|A_0|}.
\]
Assume that priorities are supported on the unit cube on every activation
atom:
\[
\nu(A=A_0,\ \pi(z)\in[0,1]\text{ for every }z\in A_0)=\nu(A=A_0).
\]
Assume that for every vector \(t\in[0,1]^{V(G)}\),
\[
\nu(A=A_0,\ 0\le \pi(z)\le t_z\text{ for every }z)
=\nu(A=A_0)\prod_{z\in V(G)}t_z.
\]
Assume also the fixed-activation one-vertex comparison identity: for every
\(z\in A_0\),
\[
\nu\bigl(A=A_0,\ \pi(q)\in[0,1]\text{ for every }q\in A_0,\ 
\pi(q)<\pi(z)\text{ for every }q\in A_0\cap N_G(z)\bigr)
\]
equals
\[
\nu(A=A_0)\int_0^1 a^{|A_0\cap N_G(z)|}\,da .
\]
For distinct nonadjacent vertices \(u,v\), let
\[
I(A,\pi)=\{z\in A:\pi(z)>\pi(w)\text{ for every }w\in A\cap N_G(z)\}.
\]
Then
\[
\nu\{(A,\pi):u\in I(A,\pi)\text{ and }v\in I(A,\pi)\}
=
\operatorname{ofReal}\!\left[
\frac{2}{\Delta^2}\int_0^\gamma\int_x^\gamma
\left(1-\frac{x}{\Delta}\right)^{
\Delta-|N_G(u)\cap N_G(v)|}
\left(1-\frac{y}{\Delta}\right)^\Delta\,dy\,dx\right].
\]
\end{helper}
\begin{proof}
The hypotheses \(0<\gamma\le\Delta\) imply \(\Delta>0\) and
\(0<\gamma/\Delta\le 1\).  Thus the activation atom formula is a genuine
Bernoulli product law with parameter \(p=\gamma/\Delta\).  The lower-rectangle
law on an atom \(A=A_0\), together with the unit-cube support, identifies the
conditional priority law with product Lebesgue measure on the finite cube
\([0,1]^{V(G)}\): lower rectangles
\(\prod_z[0,t_z]\), \(0\le t_z\le1\), form an intersection-stable generating
class for the finite product Borel algebra, and both finite measures agree on
that class and on the whole cube.  Consequently every finite Boolean
combination of coordinate intervals inside the cube has its ordinary product
volume, and disjoint finite unions of such cells have the sum of those product
volumes.  This is the lower-rectangle-to-interval conversion used below.

First fix the activation and priority coordinates of \(u\) and \(v\), and
write
\[
x=\gamma(1-\pi(u)),\qquad y=\gamma(1-\pi(v)).
\]
The affine map \((\pi(u),\pi(v))\mapsto(x,y)\) sends the unit square to
\([0,\gamma]^2\) and has Jacobian \(\gamma^2\) in the inverse direction.
Since the two priority coordinates are independent uniforms, the
\((x,y)\)-density is \(\gamma^{-2}\,dx\,dy\).  The activations of \(u\) and
\(v\) are independent of these priorities and contribute \(p^2\).  Therefore
after the two mandatory activations are included, the base density on the
\((x,y)\)-square is
\[
p^2\gamma^{-2}\,dx\,dy=\Delta^{-2}\,dx\,dy .
\]
The diagonal \(x=y\) has two-dimensional product measure zero, so the square
is the disjoint union up to a null set of the two ordered chambers
\(x\le y\) and \(y\le x\).

We compute the chamber \(0\le x\le y\le\gamma\), equivalently
\(\pi(u)\ge\pi(v)\).  A vertex \(w\notin\{u,v\}\) can affect the event
\(\{u,v\subseteq I(A,\pi)\}\) only if it is adjacent to \(u\) or to \(v\).
If \(w\in N_G(u)\setminus N_G(v)\), then \(w\) defeats \(u\) exactly when
\(w\) is activated and \(\pi(w)>\pi(u)=1-x/\gamma\).  By the activation atom
law and the interval consequence of the lower-rectangle law, this has
probability
\[
p\cdot |\,(1-x/\gamma,1]\,|=(\gamma/\Delta)(x/\gamma)=x/\Delta .
\]
Therefore such a vertex contributes the avoidance factor \(1-x/\Delta\).

If \(w\in N_G(v)\), including the common-neighbor case
\(w\in N_G(u)\cap N_G(v)\), assign \(w\) to the \(v\)-test on this chamber.
It defeats \(v\) exactly when it is activated and
\(\pi(w)>\pi(v)=1-y/\gamma\), which has probability \(y/\Delta\).  Hence it
contributes the factor \(1-y/\Delta\).  For a common neighbor this single
condition is the stronger one on the chamber: \(x\le y\) gives
\(\pi(u)\ge\pi(v)\), so \(\pi(w)\le\pi(v)\) automatically implies
\(\pi(w)\le\pi(u)\).  Thus assigning common neighbors to the \(v\)-side
prevents them from defeating either of the two surviving vertices and counts
each common neighbor once.

The finitely many activation indicators and priority coordinates for
different vertices are independent by the Bernoulli atom formula and by the
finite product priority law.  Hence the avoidance factors multiply.  Since
\(G\) is \(\Delta\)-regular, \(u\) has exactly \(\Delta\) neighbors, of which
\(|N_G(u)\cap N_G(v)|\) are common with \(v\); the \(u\)-only class therefore
has size \(\Delta-|N_G(u)\cap N_G(v)|\).  The whole neighbor set of \(v\) has
size \(\Delta\).  The chamber contribution is consequently
\[
\Delta^{-2}\int_0^\gamma\int_x^\gamma
\left(1-\frac{x}{\Delta}\right)^{
\Delta-|N_G(u)\cap N_G(v)|}
\left(1-\frac{y}{\Delta}\right)^\Delta\,dy\,dx .
\]

On the opposite chamber \(y\le x\), the same argument with \(u\) and \(v\)
interchanged gives the same numerical contribution after renaming the two
integration variables.  Adding the two ordered chambers and ignoring the null
diagonal gives the displayed factor \(2/\Delta^2\).  The integrand is
nonnegative on the chamber because \(0\le x\le y\le\gamma\le\Delta\), so the
real value inside \(\operatorname{ofReal}\) is nonnegative and represents the
\(\nu\)-mass of the survival event.
\end{proof}
